\subsection{Wybrane fragmenty kodu}
Poniższe listingi są fragmentami rzeczywistej implementacji. Pominięto
kod otaczający, aby wyeksponować najważniejsze mechanizmy.

\begin{lstlisting}[language=Java,caption={BuyerAgent.java: utworzenie żądania rezerwacji (wiersze 397--408)}]
private void sendReservation() {
    ACLMessage request = new ACLMessage(ACLMessage.REQUEST);
    request.addReceiver(selectedSeller);
    request.setOntology(MarketplaceProtocol.RESERVATION);
    request.setConversationId(conversationId);
    expectedReplyToken = UUID.randomUUID().toString();
    request.setReplyWith(expectedReplyToken);
    request.setContent(codec.encode(new MarketplaceProtocol.ReservationRequest(
            selectedOffer.car().id(), getLocalName(), agreedTotalPln)));
    send(request);
    state = WAITING_FOR_RESERVATION;
}
\end{lstlisting}

\begin{lstlisting}[language=Java,caption={SellerInventory.java: atomowa kontrola i zapis rezerwacji (wiersze 45--67)}]
public synchronized ReservationResult reserve(
        String carId,
        String buyerName,
        long agreedTotalPln,
        long durationMillis) {
    releaseExpiredReservations();
    InventoryItem item = items.get(carId);
    if (item == null) {
        return new ReservationResult(false, carId, "UNKNOWN_CAR");
    }
    if (item.status != CarStatus.AVAILABLE) {
        return new ReservationResult(false, carId, "CAR_UNAVAILABLE");
    }
    if (agreedTotalPln < NegotiationService.minimumTotal(item.car)
            || agreedTotalPln > item.car.totalCostPln()) {
        return new ReservationResult(false, carId, "INVALID_NEGOTIATED_PRICE");
    }
    item.status = CarStatus.RESERVED;
    item.reservedBy = buyerName;
    item.reservedTotalCostPln = agreedTotalPln;
    item.reservedUntilEpochMillis = clock.millis() + durationMillis;
    return new ReservationResult(true, carId, "RESERVED");
}
\end{lstlisting}

\begin{lstlisting}[language=Java,caption={SellerInventory.java: kontrola właściciela przed sprzedażą (wiersze 69--76)}]
public synchronized Optional<Transaction> purchase(String carId, String buyerName) {
    releaseExpiredReservations();
    InventoryItem item = items.get(carId);
    if (item == null || item.status != CarStatus.RESERVED || !buyerName.equals(item.reservedBy)) {
        return Optional.empty();
    }
    item.status = CarStatus.SOLD;
    long transactionTotalCostPln = item.reservedTotalCostPln;
\end{lstlisting}

\begin{lstlisting}[language=Java,caption={RecommendationService.java: normalizacja kryteriów kosztowych i korzyści (wiersze 79--85)}]
private double lowerIsBetter(double value, Range range) {
    return range.isConstant() ? 1.0 : (range.maximum() - value) / (range.maximum() - range.minimum());
}

private double higherIsBetter(double value, Range range) {
    return range.isConstant() ? 1.0 : (value - range.minimum()) / (range.maximum() - range.minimum());
}
\end{lstlisting}
Stała wartość kryterium daje wszystkim ofertom ocenę 1, dzięki czemu nie
dochodzi do dzielenia przez zero. Normalizacja jest względna wobec
badanego zbioru; wyników MCDA z różnych zbiorów nie należy traktować
jako bezwzględnej miary jakości pojazdu.
